% =============================================================================
% Template: proof-block.tex
% Skill:    iacr-math-prose
%
% Use for: a stand-alone proof body, when the theorem is stated elsewhere
%          (e.g., the theorem is in section 3 and the proof is deferred to
%          section 6 or appendix A).
%
% Invariants (do not delete):
%   - Begin with "Proof of \cref{thm:...}" so the reader can locate the claim.
%   - Use \paragraph{Name.} for sub-blocks; never \\ \\ as a separator.
%   - End with \qed (auto-inserted by amsthm if the proof ends in plain text;
%     manual if it ends in a displayed equation).
%   - Cite every external lemma by name and source.
%   - State every hardness assumption invoked explicitly.
%
% See:
%   - references/proof-style-game-based.md
%   - references/proof-style-uc.md
%   - references/proof-style-concrete.md
%   - references/common-rejection-reasons.md (R23, R29 in particular)
% =============================================================================

\begin{proof}[Proof of \cref{thm:\TODO{theorem-label}}]
\TODO{One-sentence setup paragraph: restate the assumption, the adversary class,
and the target bound. This serves as a recap so the proof body does not require
the reader to flip back to the theorem statement.}

\paragraph{Setup of the reduction.}
\TODO{Describe the reduction construction. Given an adversary $A$ against the
target notion, construct an adversary $B$ against the underlying hardness
assumption. State $B$'s inputs and what $B$ outputs.}

\paragraph{Simulation of $A$'s view.}
\TODO{Describe how $B$ simulates $A$'s oracle responses. For each oracle:
- Random oracle queries: lazy table with DST-prefixed keys.
- Signing/decryption queries: programmed using $B$'s inputs.
- Other queries: state explicitly.
State that the simulation is statistically (or perfectly, or computationally)
indistinguishable from the real game, with bound \TODO{$\varepsilon_\textsf{sim}$}.}

\paragraph{Extracting from $A$'s output.}
\TODO{Describe how $B$ uses $A$'s output (forgery, distinguishing bit, …) to
produce the assumption-breaker's output. State the success probability:
``$B$ succeeds whenever $A$ succeeds and event $E$ holds, where $\Pr[E] \geq
\TODO{...}$.''}

\paragraph{Bound.}
\TODO{Concrete bound, e.g.}
\[
  \TODO{\Adv}^{\TODO{\textsf{ASSUMPTION}}}_{B}(\secparam)
  \;\geq\;
  \frac{1}{\TODO{q_H}} \cdot
  \TODO{\Adv}^{\TODO{\textsf{NOTION}}}_{A}(\secparam)
  \;-\; \TODO{\varepsilon_\textsf{sim}}.
\]
Rearranging,
\[
  \TODO{\Adv}^{\TODO{\textsf{NOTION}}}_{A}(\secparam)
  \;\leq\;
  \TODO{q_H} \cdot
  \TODO{\Adv}^{\TODO{\textsf{ASSUMPTION}}}_{B}(\secparam)
  \;+\; \TODO{q_H \cdot \varepsilon_\textsf{sim}}.
\]

\paragraph{Runtime.}
$B$ runs in time
$\TODO{t_A + O(q_H \cdot T_\textsf{rand-oracle} + q_S \cdot T_\textsf{sign})}$,
where $\TODO{T_\textsf{rand-oracle}}$ is the cost of one random-oracle table
lookup and $\TODO{T_\textsf{sign}}$ is the cost of one signing simulation.

This completes the proof. \qed
\end{proof}
